% =====================================================================
%  §7  Literature verdict
% =====================================================================
\section{Literature verdict: the shape is known, the theorems are zero}
\label{sec:lit}

This section separates two things on the literature side: \textbf{the shape} ($m=C\Lambda$) has
numerical support and two-dimensional precedent; \textbf{$a$-uniform theorems} are zero. Both hold at
once, and that is precisely the external version of the verdict of \S\ref{sec:position}.

\subsection{The shape: $m/\Lm=6.3(3)$ at the continuum limit, three independent routes agreeing}

\begin{table}[h]
\centering\footnotesize
\setlength{\tabcolsep}{5pt}
\renewcommand{\arraystretch}{1.12}
\begin{tabularx}{\textwidth}{@{}l X l@{}}
\toprule
\textbf{Route} & \textbf{Data} & \textbf{Result}\\
\midrule
$m/\sqrt{\sigma}\div\Lm/\sqrt{\sigma}$ & AT2021 (\texttt{arXiv:2106.00364}, JHEP12(2021)082) Tab.34 $\div$ Tab.15:
$\mathrm{SU}(3)$ $3.405(21)/0.5424(13)$; $\mathrm{SU}(2)$ $3.781(23)/0.5806(21)$ & $6.28(2)(21)$ / $6.51(2)(64)$\\
$r_0$ route ($r_0\Lm=0.629(4)$) & MP1999 \texttt{hep-lat/9901004} \S V: $r_0m=4.21(12)$; Chen2006 \texttt{hep-lat/0510074}
Tab.XXI: $4.16(11)(4)$ & $6.69(20)(23)$ / $6.61(18)(24)$\\
GeV division & AT2021 Tab.37: $1.653(26)$ GeV $\div$ eq.22: $0.263(4)(9)$ GeV & $6.28(14)(22)$\\
\bottomrule
\end{tabularx}
\caption{$m_{0^{++}}/\Lm^{(N_f=0)}=6.3(3)$: three mutually independent routes, with $\mathrm{SU}(2)$ and
$\mathrm{SU}(3)$ agreeing.}
\label{tab:c}
\end{table}

\noindent
Together with the uniqueness \eqref{eq:unique} of \S\ref{sec:position} the target is now: \textbf{prove
$m\ge c\Lambda$ with $c>0$ a pure number}; the value is known, $C=6.3(3)$ (about $5\%$ precision). What
is missing is the proof, not the number. \LE
\src{\code{L2-GLUEBALL-LAMBDA.md}, as reported in \code{LOOP-RUNBOOK.md} \S10.16--\S10.17 and
\code{LIT-VERDICT.md} \S6; the scheme ratio $\Lm/\Lambda_I=2.633$ is from \code{L2-GLUEBALL-LAMBDA.md}.}

\subsection{$a$-uniform theorems: zero (and not ``not found'').}

\begin{itemize}
\item \textbf{The only rigorous mass gap on the lattice}: Kupiainen, \emph{CMP} \textbf{73} (1980) 273
      (\texttt{doi:10.1007/BF01197703}) proves $m>0$ and a $1/n$ expansion, \textbf{with no shape}. \LT
\item \textbf{Constructive 4D gauge theory}: Magnen--Rivasseau--Sénéor, \emph{CMP} \textbf{155} (1993),
      title verbatim ``\emph{with an infrared cutoff}'' --- the infrared is cut off, no $m$; the Balaban
      series (\emph{CMP} 102/109/116) proves UV stability and a small-field effective action, no $m$. \LT
\item \textbf{The only $m$ bound is outside that line}, Osterwalder--Seiler 1978: $ma\ge c(\beta)$, so
      $a\to0$ gives $m\ge c(\beta)/a\to\infty$ --- \textbf{the direction is reversed}. \LT
\item \textbf{Precedents for the shape exist only in 2D and in SUSY}: seven of them (O(3)/O(4),
      O($N$) $\sigma$ models, SUSY $\mathbb{CP}^{N-1}$, SUSY O($N$), etc.), all relying on
      integrability or SUSY; 4D has neither. \LT
\end{itemize}
\src{\code{LIT-VERDICT.md} \S1--\S4 (every DOI checked); \code{L1-UNIFORM-A.md},
\code{L3-TRANSMUTATION.md}, \code{L5-CLASSICAL.md}.}

\begin{table}[h]
\centering\footnotesize
\setlength{\tabcolsep}{5pt}
\renewcommand{\arraystretch}{1.12}
\begin{tabularx}{\textwidth}{@{}>{\raggedright\arraybackslash}p{3.5cm} >{\raggedright\arraybackslash}p{3.6cm} X@{}}
\toprule
\textbf{Source} & \textbf{Verbatim} & \textbf{Verdict}\\
\midrule
\texttt{arXiv:2606.19362}\newline(Faizal--Shabir, \emph{Fortschr.\ Phys.} \textbf{74} (2026) e70097)
 & ``transfer operator with a uniform gap \dots\ endure in the continuum''
 & \textbf{downgraded to a claim}: no constant, no display number; the core is \code{Hypothesis 1};
 uniform-in-$a$ comes from a \emph{postulated admissible class}; contraction to the Gaussian fixed
 point; no numerical bound on $\Delta_{\mathrm{phys}}$\\
\texttt{arXiv:2506.00284}\newline(Jacobsen) & ``no remaining gap'' & \textbf{excluded}: arXiv marks it
 withdrawn (``does not meet arXiv's research content quality standards'')\\
\texttt{arXiv:2307.00788}\newline(Lim) & gap $=$ Casimir label & dimensionless, no scaling, not a lattice
 $a\to0$ limit\\
\bottomrule
\end{tabularx}
\caption{The three preprints claiming to settle 4D: all die at the same step.}
\label{tab:preprints}
\end{table}

\noindent
The furthest point reached by the other nine is ``fixed-lattice-spacing strong coupling / 't Hooft
region'', and the field's own verdict is already in print: \texttt{arXiv:2611.11534} (2026-11-08)
verbatim ``\emph{A nontrivial four-dimensional spatial-cutoff construction and a positive spectral
bound along it are not established.}''; Zenodo 23783328 verbatim ``\emph{the composition across
scales, the large fields and the crossover remain open. The cube estimate and the existence of the
continuum limit are not proved, so the continuum Yang--Mills existence and mass-gap problem remains
open.}'' \LT
\src{\code{LIT-VERDICT.md} \S4--\S5 and \S7; \code{L4-PREPRINTS.md} (twelve IDs, each opened).}

\begin{keybox}
\textbf{The common gap of the three claims is the object of \S\ref{sec:five}.} Each of them, in
transporting ``strong-coupling gap $\to$ weak-coupling continuum window'', turns a hypothesis into a
theorem; and that step is exactly the rate uniformisation of $C5$, i.e.\ $d5/(CP)=p=1$. This is the
strongest external confirmation in this document: \textbf{every attempt fails at the same step, and we
know the step's name.} \LT
\end{keybox}
\src{\code{LOOP-RUNBOOK.md} \S10.18 (lines 1006--1019); \code{LIT-VERDICT.md} \S7, conclusion 4.}

\subsection{A premise that must be attached to $6.3(3)$: it is conditional on the scaling form}

The value $6.3(3)$ is obtained by extrapolating $m/\sqrt{\sigma}$ and $r_0m$ measured at
\emph{accessible} $\beta$ to the continuum limit, so it depends on the scaling form $m=C\Lambda$ itself.
At accessible $\beta$ the measured logarithmic slope is only $0.53\,(\mathrm{SU}(3))$ and
$0.21\,(\mathrm{SU}(2))$ of the asymptotic value, and for $\mathrm{SU}(2)$ the $N_\tau=4$ data even
have the opposite sign; reaching $90\%$ of the asymptotic slope needs $\beta_{90}\approx27$ (where
$\xi/a\sim10^{11}$), which is inaccessible. \LE
\src{\code{C2-MEASURED-SLOPE.md}, as reported in \code{LOOP-RUNBOOK.md} \S10.20 (including the
extrapolation formula for $\beta_{90}$ and the $1\sigma$ interval $[6.3,58]$).}

\begin{keybox}
\textbf{Hence $6.3(3)$ is ``compatible with the scaling form'', not ``established by it''.} The
numerical route \textbf{cannot independently confirm} the shape --- it takes the very thing to be proved
as the hypothesis of its extrapolation. \LN
\end{keybox}
\src{\code{LOOP-RUNBOOK.md} \S10.20, the honest-premise paragraph (lines 1064--1067).}

\noindent
This is the same statement as \S\ref{sec:position}, seen from the literature side: the shape is
numerically confirmed to $5\%$ precision, while the route to it presupposes $m>0$. That is what the
title means by ``a consequence, not a route''.

